\section{Conclusion}
\label{sec:conclusion}

This study defines and evaluates a charge-self-consistent extension of MACE for energetic molecular crystals. By coupling an equivariant short-range model to a differentiable charge-equilibration module within one periodic total-energy framework, CSC-MACE isolates explicit self-consistency as a concrete and testable modeling increment.

Our evaluation across the public reference system, two public surrogate systems, and the final CL-20/HMX benchmark shows that charge-aware effects are physically active and numerically consequential, but not uniformly beneficial. On the final energetic-material target, the static hierarchy is now unambiguous under four held-out views: the original contiguous slice, the cleaner material-balanced split-v2 evaluation, the order-recovered block-holdout split-v3 evaluation, and the stronger recovered temperature-gap split-v4 evaluation. In all four cases, \texttt{self\_consistent} yields the lowest energy error, followed by \texttt{static\_charge}, with \texttt{short\_range} the weakest.

The associated QEq diagnostics confirm that the electrostatic branches are measurably active rather than inert, with numerically negligible neutrality residuals and slightly stronger charge/field response in the self-consistent family. The final static gain is therefore not reducible to the original surrogate-system choice or to the original HMX-only contiguous split.

At the same time, the real-target nonequilibrium trajectories expose a hard dynamical boundary. Across stress, threshold-aware, reaction-progress, dense-ensemble, focused key-bond analyses, and train-matched fragile-frame comparisons, the dominant variation remains frame selection rather than model family, and in the final focused comparison the frame-driven spread exceeds the family spread by roughly $8\times$ for the first integrity-crossing step and $184\times$ for the first drift-threshold step.

The results presented here therefore argue for a layered validation strategy in which static error, early structural onset, fragment integrity, and post-onset stability are treated as distinct targets rather than interchangeable summaries. For shock-driven molecular crystals, that multidimensional view is essential. The present CSC-MACE variants improve static energetics on the real target, but they do not induce a corresponding global family-level separation in the tested nonequilibrium observables.

The strongest positive dynamic signal we now find on CL-20/HMX is a narrow local endpoint-survival window rather than a global hierarchy. In the dense fragile-window analysis, train-matched self-consistency improves endpoint structural survival on frames $64$ and $72$, with the clearest local gain on frame $72$, where final molecular integrity rises by about $0.058$ relative to \texttt{static\_charge}, the first integrity crossing is delayed by about $195$ steps, and final nitro survival increases slightly as well. That window does not extend to frame $56$ and does not establish robust dynamic superiority over the target as a whole.

Even stronger adaptive interventions on the same fragile-frame ladder generally worsen early onset while increasing rollout cost, showing that more adaptive self-consistency alone does not create a new dynamical hierarchy. The resulting priority is clear: the next challenge is to stabilize the electrostatic gain so that it propagates into dynamically relevant shock observables without disproportionate computational overhead.

The present data boundary is also clear: the CL-20/HMX release already supports a real force-labeled energetic-material benchmark, but it does not yet retain the reviewer-complete thermodynamic and labeling provenance needed for a final pressure- or trajectory-history-resolved mechanistic decomposition.
